\section{O Teorema do Confronto}\index{Teorema do Confronto}\index{sanduíche, teorema do|see{Teorema do Confronto}}


\begin{proposition}Sejam
    \begin{itemize}
        \item $A\subseteq \mathbb R^n$ um conjunto,
        \item $f, g: A \to \mathbb R$ funções,
        \item $p \in \mathbb R^n$ um ponto de acumulação de $A$,
        \item $L=\lim_{x\to p} f(x)$, e,
        \item $M=\lim_{x\to p} g(x)$.
    \end{itemize}

    Então:
    \begin{enumerate}[label=(\alph*)]
        \item $\displaystyle \lim_{x\to p}(f+g)(x) = L + M$,
        \item $\displaystyle \lim_{x\to p}(f-g)(x) = L - M$, e 
        \item $\displaystyle \lim_{x\to p}(fg)(x) = LM$.
    \end{enumerate}
\end{proposition}

\begin{proposition}Sejam
    \begin{itemize}
        \item $A\subseteq \mathbb R^n$ um conjunto,
        \item $f, g: A \to \mathbb R$ funções,
        \item $p \in A$.
    \end{itemize}

    Se $f$ e $g$ são contínuas em $p$, então:
    \begin{enumerate}[label=(\alph*)]
        \item $f+g$ é contínua em $p$,
        \item $f-g$ é contínua em $p$, e
        \item $fg$ é contínua em $p$.
    \end{enumerate}

    Consequentemente, se $f$ e $g$ são contínuas, então $f+g$, $f-g$ e $fg$ são contínuas.
\end{proposition}

\begin{proposition}Sejam
    \begin{itemize}
        \item $A\subseteq \mathbb R^n$ um conjunto,
        \item $f, g: A \to \mathbb R$ funções com $g(x)\neq 0$ para todo $x \in A$,
        \item $p \in \mathbb R^n$ um ponto de acumulação de $A$,
        \item $L=\lim_{x\to p} f(x)$, e,
        \item $M=\lim_{x\to p} g(x)$.
    \end{itemize}

    Então:
    \[
        \lim_{x\to p}\frac{f(x)}{g(x)} = \frac{L}{M}.
    \]
\end{proposition}

\begin{proposition}Sejam
    \begin{itemize}
        \item $A\subseteq \mathbb R^n$ um conjunto,
        \item $f, g: A \to \mathbb R$ funções com $g(x)\neq 0$ para todo $x \in A$,
        \item $p\in A$.
    \end{itemize}

    Se $f$ e $g$ são contínuas em $p$, então $\displaystyle\frac{f}{g}$ é contínua em $p$.

    Consequentemente, se $f$ e $g$ são contínuas, então $\displaystyle\frac{f}{g}$ é contínua.
\end{proposition}


\begin{proposition}\index{limite!Teorema do Confronto}
    Sejam $f, g, h: A \to \mathbb R$ uma função, onde $A\subseteq \mathbb R^m$.

    Sejam $p \in A$ um ponto de acumulação de $A$ e $L \in \mathbb R$.

    Suponha que existe uma bola aberta $B$ em torno de $p$ tal que, para todo $x \in B\cap A\setminus\{p\}$, temos que $f(x) \leq g(x) \leq h(x)$.

    Nessas hipóteses, se $\lim_{x\to p} f(x) = \lim_{x\to p} h(x) = L$, então $\lim_{x\to p} g(x) = L$.
\end{proposition}

\begin{proposition}\index{limite!Teorema do Confronto}
    Sejam
    \begin{itemize}
        \item $A\subseteq \mathbb R^n$ um conjunto,
        \item $f, g, h: A \to \mathbb R$ funções,
        \item $p \in \mathbb R^n$ um ponto de acumulação de $A$, e
        \item $L \in \mathbb R$.
    \end{itemize}

    Suponha que existe uma bola aberta $B$ em torno de $p$ tal que, para todo $x \in B\cap A\setminus\{p\}$, temos
    \[
        f(x) \leq g(x) \leq h(x).
    \]

    Então:
    \[
        \text{se } \lim_{x\to p} f(x) = \lim_{x\to p} h(x) = L, \text{ então } \lim_{x\to p} g(x) = L.
    \]
\end{proposition}

\begin{proof}
    Seja dado $\epsilon > 0$.
    
    Como $\lim_{x\to p} f(x) = L$, existe $\delta_f>0$ tal que, para todo $x \in A\setminus\{p\}$, se $d(x, p) < \delta_f$, então $|f(x) - L| < \epsilon$.

    Analogamente, como $\lim_{x\to p} h(x) = L$, existe $\delta_h>0$ tal que, para todo $x \in A\setminus\{p\}$, se $d(x, p) < \delta_h$, então $|h(x) - L| < \epsilon$.

    Por fim, existe $\delta_g>0$ tal que, para todo $x \in B(p, \delta_g)$, se $x \in A\setminus\{p\}$, então $f(x) \leq g(x) \leq h(x)$.

    Seja $\delta = \min\{\delta_f, \delta_h, \delta_g\}$.

    Seja $x \in A\setminus\{p\}$ tal que $d(x, p) < \delta$, então $-\epsilon<f(x) - L \leq g(x) - L \leq h(x) - L < \epsilon$, logo, $|g(x) - L| < \epsilon$, ou seja, $d(g(x), L) < \epsilon$.
\end{proof}
\begin{corollary}
    Sejam
    \begin{itemize}
        \item $A\subseteq \mathbb R^n$ um conjunto,
        \item $f, g$ funções,
        \item $p \in \mathbb R^n$ um ponto de acumulação de $A$, e
        \item $L \in \mathbb R$.
    \end{itemize}

    Suponha que:
    \begin{itemize}
        \item existe uma bola aberta $B$ em torno de $p$ tal que $g|_B$ é limitada, ou seja, tal que existe $M>0$ tal que, para todo $x \in B\cap A\setminus\{p\}$, temos $|g(x)| < M$, e
        \item $\lim_{x\to p} f(x) = 0$.
     \end{itemize}

     Então $\lim_{x\to p} f(x)g(x) = 0$.
\end{corollary}

\begin{proposition}Sejam
    \begin{itemize}
        \item $I\subseteq \mathbb R$ um intervalo aberto
        \item $A\subseteq \mathbb R^n$ um conjunto,
        \item $a \in \mathbb R^n$ um ponto de acumulação de $A$,
        \item $\gamma: I \to \mathbb R^n$ uma função
        \item $t_0 \in I$ com $\gamma(t_0) = a$
    \end{itemize}

    Assuma que:
    \begin{itemize}
        \item $\gamma(t)\neq a$ para todo $t \in I\setminus\{t_0\}$,
        \item $\lim_{t\to t_0} \gamma(t) = a$, e
        \item $\lim_{t\to t_0} f(\gamma(t)) = L$.
    \end{itemize}

    Então $\lim_{t\to t_0} f(\gamma(t)) = L$.
\end{proposition}